\section{官方幻灯片证据链：Self-Play 的适用条件与外推边界}

本节按课堂顺序保留 Noam Brown 的完整论证。主线不是“多放几个 Agent 就更强”，而是先问博弈结构：二人零和、完美信息、不完美信息和多人 general-sum 对目标函数与算法有完全不同的要求。最后再回到 LLM cooperation，检查 consensus、并行推理和 diversity 是否真的抵消延迟与成本。

\subsection{AlphaGo 三步：Human Data、Inference Compute、Self-Play}

slides 2--6 把 AlphaGo 路线概括为三步：用高质量人类数据预训练；在推理时用搜索投入更多 compute；通过 self-play 递归改进。LLM 已广泛采用前两步，第三步是否可直接复制仍是问号。self-play 成功需要环境规则稳定、reward 清楚、对手分布由系统控制，并且算法能够在对应博弈类别中收敛。

\lecturefigure{slides-images/slide-002.png}{官方 slide 2：游戏中的 Self Play}
\lecturefigure{slides-images/slide-003.png}{官方 slide 3：AlphaGo 三步类比}
\lecturefigure{slides-images/slide-004.png}{官方 slide 4：Step 1 高质量人类棋谱预训练}
\lecturefigure{slides-images/slide-005.png}{官方 slide 5：Step 2 大规模 inference search}
\lecturefigure{slides-images/slide-006.png}{官方 slide 6：Step 3 recursive self-improvement 的问号}

\begin{knowledgebox}{读图：三步解决不同问题}
人类数据提供合理策略先验，搜索在当前模型上扩大决策深度，self-play 改变训练分布并发现超越人类的新策略。把三者混成“更多 compute”会漏掉数据来源和对手分布。LLM tool environment 还存在开放规则、不可验证 reward 和真实用户，因此第三步最难复制。
\end{knowledgebox}

\subsection{Minimax 与 Exploitability：稳健最优不等于平均收益最高}

slides 7--17 用扑克选手问题引出 minimax equilibrium。面对任意对手，minimax 策略保证最坏情况；exploitability 衡量对手用 best response 能额外赢多少。slides 9--13 逐步展示混合策略使 exploitability 从 1 降到 0；slides 14--15 说明 sound self-play 需要初始策略覆盖和收敛算法；完美信息二人零和可近似单 Agent RL，但仍可能受 adversarial attack。

\lecturefigure{slides-images/slide-007.png}{官方 slide 7：平均赢得多与被最强对手难以利用的冲突}
\lecturefigure{slides-images/slide-008.png}{官方 slide 8：Minimax equilibrium 作为另一种 better player 定义}
\lecturefigure{slides-images/slide-009.png}{官方 slide 9：Minimax 与 exploitability 初始示例}
\lecturefigure{slides-images/slide-010.png}{官方 slide 10：对手 best response 与混合概率}
\lecturefigure{slides-images/slide-011.png}{官方 slide 11：调整混合策略降低 exploitability}
\lecturefigure{slides-images/slide-012.png}{官方 slide 12：Our exploitability 降为 0}
\lecturefigure{slides-images/slide-013.png}{官方 slide 13：Mastering Poker 的 minimax 解释}
\lecturefigure{slides-images/slide-014.png}{官方 slide 14：二人零和 self-play 的改善与收敛}
\lecturefigure{slides-images/slide-015.png}{官方 slide 15：不同游戏中的 self-play}
\lecturefigure{slides-images/slide-016.png}{官方 slide 16：完美信息二人零和与 single-agent RL}
\lecturefigure{slides-images/slide-017.png}{官方 slide 17：完美信息策略仍可能遭遇 adversarial attack}

\begin{importantbox}{Exploitability 定义}
若策略为 $\pi$，对手 best response 为 $BR(\pi)$，exploitability 可理解为对手相对 equilibrium value 能多获得的收益。它评价最坏对手下的漏洞，不等于对真实 population 的平均回报。低 exploitability 适合 adversarial setting，却可能放弃对弱对手的额外收益。
\end{importantbox}

\begin{warningbox}{公式的适用域}
Minimax 的强保证依赖二人、零和和正确建模。多人或非零和中，联盟、沟通和不同效用会破坏唯一解释；把 exploitability 当通用 Agent 指标，会忽略合作价值与人类分布。\textbf{课堂提示：}讲者反复比较两种“better poker player”，就是在提醒目标定义先于算法。
\end{warningbox}

\subsection{不完美信息：PPO 循环与 Regret 家族}

slides 18--30 从 Rock-Paper-Scissors 展示 PPO 可能围绕策略循环，无法稳定达到混合 equilibrium。不完美信息中，同一可见状态对应不同隐藏历史，动作价值取决于整套混合概率。Fictitious Play 对对手历史平均策略做 best response；Regret Matching/Hedge 根据未选动作的 counterfactual regret 调整概率；改进算法与 Libratus 结果表明低 exploitability 可在大规模扑克成立，但 FP/RM 在 single-agent RL 可能慢或表现差。

\lecturefigure{slides-images/slide-018.png}{官方 slide 18：PPO 在 Rock-Paper-Scissors 中循环}
\lecturefigure{slides-images/slide-019.png}{官方 slide 19：RPS+ 的隐藏信息树}
\lecturefigure{slides-images/slide-020.png}{官方 slide 20：同一 observation 下不同 history}
\lecturefigure{slides-images/slide-021.png}{官方 slide 21：策略概率改变 action value}
\lecturefigure{slides-images/slide-022.png}{官方 slide 22：不完美信息中的完整混合策略}
\lecturefigure{slides-images/slide-023.png}{官方 slide 23：Fictitious Play 与历史平均 best response}
\lecturefigure{slides-images/slide-024.png}{官方 slide 24：Regret Matching、Hedge 与学习率}
\lecturefigure{slides-images/slide-025.png}{官方 slide 25：改进 regret matching 的收敛曲线}
\lecturefigure{slides-images/slide-026.png}{官方 slide 26：2017 Brains vs AI Libratus 对战}
\lecturefigure{slides-images/slide-027.png}{官方 slide 27：120,000 hands 与统计结果}
\lecturefigure{slides-images/slide-028.png}{官方 slide 28：FP/RM/Hedge 在 single-agent RL 的限制}
\lecturefigure{slides-images/slide-029.png}{官方 slide 29：Last-iterate convergence 新算法}
\lecturefigure{slides-images/slide-030.png}{官方 slide 30：二人零和 equilibrium 中 cheap talk 无效}

\begin{knowledgebox}{术语消化：Average-iterate 与 Last-iterate}
许多 regret 算法保证历史平均策略收敛，但训练最后一个 checkpoint 可能仍振荡；last-iterate convergence 要求当前策略本身趋近 equilibrium。部署只能使用具体策略，因此 last-iterate 对工程更直接。regularization 可改善稳定性，却也改变原博弈和最优响应。
\end{knowledgebox}

\textbf{老师强调：}Libratus 的证据是特定二人零和扑克中的 sound strategy，不是证明 regret 方法适合所有 RL。算法应匹配博弈结构；单 Agent 环境、多人协作或语言谈判需要不同目标和数据。

\subsection{General-Sum 与人类 Population：Support 比单一均衡更关键}

slides 31--42 是核心转折。多人非零和中，minimax 可能无意义；Ultimatum Game 的人类接受分布受文化影响；Diplomacy 包含联盟和协调。DORA 从 self-play 学到 no-press Diplomacy，但不同训练 population 交叉对战会失配，support coverage 很关键。课程提出把人类当环境：先模仿人类数据，再用 self-play/RL 扩展。CICERO 与真实玩家、Hanabi cross-play 结果说明评估必须覆盖人类和多 population。

\lecturefigure{slides-images/slide-031.png}{官方 slide 31：General-sum 中 minimax 失去通用解释}
\lecturefigure{slides-images/slide-032.png}{官方 slide 32：Ultimatum Game 与人类接受分布}
\lecturefigure{slides-images/slide-033.png}{官方 slide 33：Diplomacy 的联盟、沟通与多人目标}
\lecturefigure{slides-images/slide-034.png}{官方 slide 34：1v1、2v1 等 support coverage}
\lecturefigure{slides-images/slide-035.png}{官方 slide 35：DORA no-press Diplomacy self-play}
\lecturefigure{slides-images/slide-036.png}{官方 slide 36：Minimax 与 population best response 再比较}
\lecturefigure{slides-images/slide-037.png}{官方 slide 37：非零和中 minimax 不具唯一意义}
\lecturefigure{slides-images/slide-038.png}{官方 slide 38：Treat humans as part of environment}
\lecturefigure{slides-images/slide-039.png}{官方 slide 39：No-Press Diplomacy population results}
\lecturefigure{slides-images/slide-040.png}{官方 slide 40：CICERO 与匿名在线人类玩家}
\lecturefigure{slides-images/slide-041.png}{官方 slide 41：Hanabi 与不同 human/agent partners}
\lecturefigure{slides-images/slide-042.png}{官方 slide 42：二人零和是特殊情况}

\begin{knowledgebox}{Population Best Response}
目标不再是抵抗任意对手，而是对目标 population 分布最大化期望收益。它需要数据描述会遇到谁，并在 subgroup 上报告表现。若训练 population 缺少某种文化、策略或工具生态，Agent 可能对主流高分却系统性伤害少数群体。
\end{knowledgebox}

\begin{warningbox}{Cross-play 是多 Agent 的基本验收}
同一训练池内 self-play 高分可能来自共适应暗号，换一个 partner 就失败。应做不同 seed、架构、训练 population 与人类的 cross-play，并分析最差群体。CICERO/Hanabi 图支持 population data 的必要性，不意味着模仿人类即可解决所有合作。
\end{warningbox}

\subsection{LLM Multi-Agent：并行、Consensus、Diversity 与 Scaffolds}

slides 44--51 将结论带回 LLM。o1 到 o3 的 inference compute 显示 test-time scaling；CoT 天生串行，多 Agent 并行可降低关键路径或做 Best-of-N/consensus，但通信和 judge 增加 latency。相同模型副本高度相关，diversity 才可能覆盖不同错误；routing 选择合适 specialist 已是一种多 Agent。现有 collaboration scaffold 可以编排，但行业文章也提醒链路复杂度可能超过收益。

\lecturefigure{slides-images/slide-044.png}{官方 slide 44：OpenAI o1 inference compute scaling}
\lecturefigure{slides-images/slide-045.png}{官方 slide 45：o1 到 o3 的能力/计算扩展}
\lecturefigure{slides-images/slide-047.png}{官方 slide 47：CoT 串行瓶颈与多 Agent latency}
\lecturefigure{slides-images/slide-048.png}{官方 slide 48：Best-of-N/consensus 的质量与延迟权衡}
\lecturefigure{slides-images/slide-049.png}{官方 slide 49：Reasoning models 的相似性与 diversity}
\lecturefigure{slides-images/slide-050.png}{官方 slide 50：当代 Multi-Agent collaboration scaffolds}
\lecturefigure{slides-images/slide-051.png}{官方 slide 51：对多 Agent 复杂度的行业反思}

\begin{knowledgebox}{读图：Consensus 的相关误差问题}
若每个 Agent 独立正确概率为 $p$，独立假设下投票可提高可靠性；但相同模型、prompt 和数据使错误高度相关，收益远低于理论。Diversity 可来自不同模型、工具、上下文、角色或搜索策略。应测 pairwise disagreement 与 error overlap，而非只增加 Agent 数。
\end{knowledgebox}

\begin{importantbox}{多 Agent 上线判据}
比较单 Agent baseline 与多 Agent 在同等 token、wall-clock 和 tool budget 下的质量；报告 p50/p95 latency、通信轮数、judge 错误和最坏 partner；只有增益超过编排与故障面成本才采用。多 Agent 是优化手段，不是架构目标。
\end{importantbox}

\subsection{读图与公式桥接：目标函数改变，最优策略也改变}

前面的 slides 从二人零和一路走到 LLM collaboration，真正变化的是“对谁最优”。在二人零和中，可写稳健目标
\[
\pi^\star=\arg\max_{\pi}\min_{\sigma}u(\pi,\sigma),
\]
其中 $\pi$ 表示我方混合策略，$\sigma$ 表示对手策略，$u(\pi,\sigma)$ 表示我方期望效用。内层 $\min$ 代表最坏对手，因此 $\pi^\star$ 追求 worst-case guarantee。Exploitability 则比较 best response 对 $\pi$ 能获得的额外收益；数值越低，策略越难被针对。

面对人类或多个 Agent population，常见目标改为
\[
\pi^\star_{\mathcal P}
=\arg\max_{\pi}\mathbb E_{\sigma\sim\mathcal P}[u(\pi,\sigma)],
\]
其中 $\mathcal P$ 表示部署时 partner/opponent 的分布。该目标可以利用常见人群，但若 $\mathcal P$ 缺少某类用户，平均最优策略可能伤害该群体。因此报告总体均值之外，还要给 subgroup、worst-group 与 cross-play。Minimax 和 population best response 不是谁更先进，而是回答不同风险偏好。

LLM consensus 还要把系统成本写进目标。设 $Q_N$ 为 $N$ 个 Agent 协作后的质量，$C_N$ 为 token/tool 成本，$L_N$ 为 latency，可比较
\[
J_N=Q_N-\lambda C_N-\mu L_N,
\]
其中 $\lambda$ 和 $\mu$ 分别是成本与延迟权重。若 $Q_N$ 的边际增益小于通信与 judge 开销，增加 Agent 数会让 $J_N$ 下降。对并行 Agent，wall-clock 由最慢分支和聚合阶段决定；对串行 debate，轮数几乎线性增加延迟。

\begin{knowledgebox}{读图：三种曲线应一起看}
质量随 Agent 数上升并不够，还要看 error overlap、latency 和 cost。若多个副本错误高度相关，投票曲线会很快饱和；若 specialist 互补，routing 可能用更少调用获得更高质量。slide 47--50 的教学结论是测量相关性和关键路径，而不是默认 consensus 总会赢。
\end{knowledgebox}

\begin{warningbox}{平均收益与公平/安全约束}
Population objective 可能牺牲少数群体；cost-adjusted objective 也可能偏向廉价但危险的策略。高风险权限、欺诈或安全违规应作为 hard constraint，不应只由 $\lambda,\mu$ 加权。工程上先满足不可违反边界，再在可行策略中优化质量、成本与延迟。
\end{warningbox}

